\section{Steering marked residues}
In a tree rooted at zero with $s$ nonroot vertices and $\ell$ marked
nonroot vertices, work modulo an odd integer $q$. We construct distinct
nonzero vertex classes up to sign which are exactly the edge classes.
Each marked vertex will have a centered representative of magnitude
$O((\ell+1)^3(s+1))$, provided $q$ exceeds a sufficiently large absolute
multiple of that bound.

At the root choose each child as a free seed $y$. At every other vertex
pair children, leaving at most one singleton. Partition each resulting
singleton path into triples followed by at most two edges. At input $p$
use seeds $y,p-y$ for siblings and $y,p-y,p-2y$ along a triple. Their
differences reproduce the seeds up to sign. Tail edges halve the parent
modulo $q$. Each occurrence is thus a seed divided by $2^t$, $0\le t\le2$.
Each seed starts at most one tail. Process a whole triple before the
families depending on its vertices.

Reserve the signed classes of $z,z/2,z/4$ for every seed $z$, requiring
its additive order to exceed 16. At any free family fewer than
$F=2000(s+3)$ parameter residues are forbidden. Indeed a linear
congruence of nonzero integer coefficient $a$ has at most $|a|$
solutions modulo $q$. Excluding orders through 16 costs at most
$3\cdot2\sum_{i=1}^{16}i$. For each old seed and new seed compare their
ratio formally by testing equality up to factors $\pm 2^j$,
$-2\le j\le2$, clearing denominators. The coefficient on $y$ has
magnitude at most 8, costing at most $240s$ across the family.
Comparisons within a family cost at most $3\cdot10\cdot16$.
If the coefficient vanishes there, the distinct linear forms in $(y,p)$
are nonproportional, so the remaining equation is $ap=0$ with
$0<|a|\le8$, impossible by the order condition. These reservations
make all used classes distinct.

A mark whose seed is halved $t$ times requests a small integer
representative of that seed divisible by $2^t$. Process the families
backwards, merging requests on a seed by taking the maximum exponent.
One requested output needs no constraint on the input. For at least
two distinct requested outputs, let $e$ be their second largest
exponent. Request that $p$ be small and divisible by $2^{e+1}$,
adding the input's tail length to this exponent at its seed.
There are at most $\ell-1$ such propagations for $\ell>0$, since each
replaces at least two active requests by at most one. An active request
of exponent $E$ represents at least $2^{(E-2)/3}$ original marks,
disjoint from those of other active requests. This holds initially.
At a propagation the largest two exponents are at least $e$, and the
new exponent is at most $e+3$. Pooling the two witness sets gives the
claim, also when merging with another request. Hence $2^E\le
4\max(1,\ell)^3$.

Put $W=16(\ell+1)^3(F+2)$. In forward order choose a singly requested
output itself in $[1,W]$ divisible by its required power of two.
Its coefficient in $y$ is a unit. With several requests the input
already has a small representative divisible by $2^{e+1}$. Choose
$y$ in $[1,W]$ or $[-W,-1]$, on the same side of zero as $p$, divisible by $2^e$, making the most constrained output
divisible by $2^E$. For $y,p-y$ the congruence is immediate. For
$p-2y$ use the class of $p/2$ modulo $2^{\max(e,E-1)}$.
Each case leaves a class modulo at most $2^E$, hence more than $F$
trials. All other requested outputs satisfy their divisibilities.
Unconstrained families avoid the original exclusions directly.
A singly requested output has magnitude at most $W$. At a multiple
request, if $|p|\le2W$, then $y,p-y,p-2y$ all have magnitude at most
$2W$ because $y$ has the sign of $p$. Halving a divisible representative
halves the integer. By induction $2W$ bounds all requested magnitudes,
and $q>8W+2F$ suffices.

\section{Two large leaf supplies}
Use the two largest demands as banks. At each remaining site reserve
one actual leaf for an odd demand or two for an even demand. The
enlarged core has at most $3n$ vertices. Set $B=10^7(k+2)^3(n+1)$. The modular forest has at most
$3n$ nonroot vertices and at most $k+1$ marks, so
$W<10^5(k+2)^3(n+1)$. Choose $P=2q$ maximal with $q$ odd and
$P\le N-2B-10$. The two banks each have at least $100B$ leaves,
so $q>8W+2F$, and $2B+10\le N-P<2B+14$. Embed the residues by
multiplication by $2$ modulo $P$, using representatives in $[1,P-1]$.
Marked representatives are within $4W<B$ of a bank. The two
banks will receive ranks $0,P$. The following reductions specify the
tree to label modulo $q$ and the ranks to restore.
\begin{itemize}
\item At distance one contract the bank edge to the root. Restore it
with length $P$.
\item At distance at least three write the path as $u,w,z,\ldots,v$.
Delete $uw$, contract $wz$ to a marked vertex, and identify $u,v$ to
form the root. The deletion separates the banks, so this quotient is
a tree. By global negation take the embedded marked residue to be
$t\in(0,B]$. Give $u,v,z,w$ ranks $0,P,t,P+t$. The two restored
lengths are $P,P+t$.
\item At distance two, if the middle $w$ has a spare neighbor $x$,
give $u,w,v,x$ ranks $0,P+t,P,t$. Their lengths are $P+t,t,P$.
If $w$ is bare and there is an auxiliary two-edge path $u,x,y$,
interchange banks if necessary and use $0,P-t,P,P+t,t$ for
$u,w,v,x,y$. The lengths are $P-t,t,P+t,P$.
In both cases replace this connected skeleton by one root-to-mark
edge, attaching outside components at the root for residue zero and
at the mark for residue $\pm t$. Solve on this tree, steering the
mark as above. Negate whole outside components attached at residue
$-t$ to restore their edge classes. If neither distance-two pattern
is available, the enlarged core is a bare double broom. No secondary
site is possible there because of its reserved leaves (a secondary
site on an end leaf would create an auxiliary path). Write its end demands as $A,C$ and label the ends and middle by
$0,1,C+2$, respectively. Give the leaves at $1$ labels
$2,\ldots,C+1$, and the leaves at $0$ labels
$C+3,\ldots,A+C+2$. The differences are exactly
$1,\ldots,A+C+2$, so this case is graceful directly.
\end{itemize}

Every secondary rank is within $B$ of zero or $P$, including restored
skeleton vertices. For ordinary edges, the unsigned residue classes of
their lengths match the ordinary vertex classes once each. In the last
two skeletons the abstract root-to-mark edge has been replaced: the
spare-neighbor pattern supplies ordinary rank and length $t$, and the
auxiliary pattern supplies both $t,P-t$ as ranks and lengths.
The special ranks and lengths $P$ and (when present) $P+t$ cancel.
Any further length above $P$ is $P+d$ with $0<d<t$, incident to the
unique vertex at $P+t$, and its unsigned class is unique. Give slot
$P+d$ to a leaf at $P$, supplying length $d$; every other unoccupied
slot above $P$ goes to zero. Each singly occupied ordinary class now
has one rank and one length. Give its other rank to zero or $P$ to
supply the missing length. The augmented tree is consistent, and
the midpoint $P/2=q$ is unused since it lies outside the subgroup of multiples of $2$ containing the modular core and all repaired ordinary classes. If the remaining demand at
$P$ is odd, give it that midpoint, supplying its own length.

The repairs spend at most $N-P+3n+1$ bank leaves: there are at most
$N-P$ high slots, fewer than $3n$ singly occupied ordinary classes,
and one midpoint. The resulting edge count $s_0$ is $O(B+n)$.
In fact $s_0\le N-P+6n+1<3B$, so each bank still has at least
$97B$ leaves available.

\subsection{Two intervals for all secondary sites}
Write $a,c$ for the residual demands at $0,P$, and $D_H,D_L$ for
the total even demands of the high and low sites, respectively.
Put $D=D_H+D_L$ and $s=s_0$, so $N-s=a+c+D$. We show that it
suffices to have $a\ge N-P+8B+10s+20$ and $c\ge2B+4s+4$, with $c$
even. Both inequalities follow from the preceding estimates.
Among the mirror indices $1,\ldots,\lfloor(P-1)/2\rfloor$ reserve
disjoint consecutive intervals of lengths $D_H/2+2B+s+2$ and
$D_L/2+B+2s+2$. They fit since $P-2\ge D+6B+6s+8$.

In the high interval visit sites of rank $h$ in decreasing order of
$d=h-P$. Starting $B$ indices after the first, scan successive $x$
and put $y=x-d$. Both $x$ and $y$ strictly increase across all trials,
including changes of site. Accept $x,P-y$ if both ranks are unused,
moving to the next site once its demand is met. The ranks sum to $h$
and reproduce their lengths in reverse order. The two halves are
disjoint, so every rejection uses a different original occupied rank.
At most $s$ trials fail, and both indices stay within the two margins.

In the low interval provisionally attach all intact mirror pairs to
$P$. At most $s$ indices are spoiled. Visit low sites in decreasing
order of $h$, scan $x$ from $B$ indices after the start, and use
$y=x-h$. Accept when $x,y$ both index originally intact pairs. Each
spoiled index can reject once as $x$ and once as $y$, so the interval
suffices. Transfer the slots $x,P-y$ to $h$. Their old lengths $P-x,y$
are exchanged. Separate transfers can use opposite members of one
mirror pair: strict increase of both index sequences ensures that
no individual slot moves twice. At most $2B+4s+4$ provisional leaves
remain at $P$, an even number. This provisional assignment describes
only the final parents and requires no extra target vertices.

Outside the intervals, intact pairs contain at least
$P-2-D-6B-8s-8\ge c$ slots by the bound on $a$. Fill the outstanding
even demand at $P$ there, and give the complement to zero. This proves
the two-supply assertion without a loss per secondary site.

\section{A rooted labeling of polynomial span}
For a rooted tree on $s\ge2$ vertices we construct consistent integer
ranks of maximum $M\le160000s^{11}$ and minimum positive rank at
least $M/(200s^6)$. Set $D=10000s^2$. At zero give its children
distinct integer ranks in $[Ds^6,2Ds^6]$ in distinct dyadic orbits
(ratios within an orbit are powers of two). Follow a largest child
at every vertex to form heavy paths. A light child in a subtree of
size $S$ has subtree size at most $S/2$.

At a path head of subtree size $S$ start with rank at least $DS^6$.
Use triples along the path, leaving at most two tail edges. From $p$
the triple uses $t,p-t,p-2t$ in that order, where
$4DJ^6\le t\le5DJ^6$ is integral and $J$ is the largest light subtree
size at the first new vertex, or one if none exists. The sets charged
by these $J$ are disjoint (charge the vertex itself when necessary),
so $\sum J\le S$ and $J\le S/2$ whenever a triple occurs.
The total boundary decrease is at most
$10D(S/2)^5S<DS^6/3$. Thus boundary and middle ranks exceed
$2DS^6/3$ and $t<p/3$. At each path vertex group its remaining children
in pairs and at most one singleton, including its heavy child if still
unassigned in the tail. At rank $z$ use $y,z-y$ with integral
$z/3\le y<z/2$, and $z/2$ for a singleton. Process the tail from
its top. Each light child of the low vertex at rank $t$ gets at least
$t/3\ge DJ^6$. Elsewhere at most two tail steps and one side step
lose a factor of at most 27 from $2DS^6/3$, leaving more than
$D(S/2)^6$. The head bound therefore holds at each recursive light
child. The ranks stay below the original head maximum and all positive
ranks are at least $D/100$.

At free choices avoid all earlier dyadic orbits and coincidences among
the new orbits. Over the trial interval each triple expression varies
by a factor less than two and any two have a strictly monotone ratio
varying by a factor less than four. This follows from $p>3\max t$ and
the displayed bounds on $t$. The sibling expressions satisfy the same
bounds. Thus at most two trials per expression and old rank, and at
most three per pair of new expressions, are forbidden. The triple
interval contains at least $D$ integers. At a sibling pair the parent
has $z\ge 4D$ or $z\ge(2/3)DS^6/9$, so the interval of width $z/6$
also suffices. In each fresh orbit forced singleton halves form only
one descending chain. The packets reproduce their ranks as differences.
Along a root path only singleton halves introduce denominators, at most
three per heavy segment, and the light steps halve subtree size.
Clearing denominators costs at most $2^{3(1+\lfloor\log_2s\rfloor)}
\le8s^3$. This gives the stated maximum and preserves the minimum ratio.

\section{Completion with one large supply}
Root at a largest site. If the second-demand hypothesis fails, the
total $L$ of secondary leaves is at most $8\cdot10^9(n+1)^5$ since $k\le n$.
This includes $k=1$ with $L=0$; zero demand is vacuous above the
threshold. Reserve one leaf at each odd secondary site. The rooted
core now has $s\le2n$ vertices. The singleton case is a star.
Otherwise scale the preceding integer ranks by $\lfloor N/M\rfloor$.
For $N\ge2M$ each secondary rank $p$ is at least $N/(400s^6)$.
The candidate pairs $\{x,p-x\}$ strictly below $p$ cover at least
$p-2$ slots. Previously occupied ranks spoil at most $2(s+L)$ slots.
Hence $p>4(s+L)+10$ suffices to install the even residual demand
at each secondary site in turn. Fill all remaining slots at zero.
For $N\ge10^{16}(n+1)^{11}$, the bounds $M\le160000s^{11}$ and
$s\le2n$ give $N\ge2M$. Also $400s^6(4(s+L)+10)<N$ by the displayed
bound on $L$. The assigned pairs and remaining root slots are leaves. This completes
the proof of Theorem~\ref{poly}.
